\subsection*{一步一步建立直观}

\paragraph{1. 选择在求导后仍保持形状的模式}

正弦函数二次求导后仍是同一个正弦形状，只多了负的频率平方。于是试设温度为一个正弦模式乘以随时间变化的振幅，偏微分方程就变成振幅的常微分方程。右图蓝线与橙线分别承担一个模式，紫线是二者相加。



\paragraph{2. 把空间弯曲转成振幅衰减}

频率为 $n$ 的模式满足振幅变化率等于 $-n^2$ 倍当前振幅，所以衰减因子为 $e^{-n^2t}$。观察同一时刻的两个振幅读数，就能把“尖细起伏消失得快”的图像翻译成定量关系。这里扩散系数取一，其他系数会相应改变时间尺度。



\paragraph{3. 叠加有效，是因为方程线性}

加大高频频率后，紫线更快接近低频模式。能够逐个演化再相加，依赖热方程的线性。若把有限个模式进一步推广成无限傅里叶级数，还需说明级数与求导、取极限之间的交换条件；两个模式的图示不能自动完成这些分析论证。



\subsection*{定义、命题与证明}

\paragraph{定义：周期热方程的一个模式}

在首尾相接的区间 $[0,2\pi]$ 上，取整数频率 $n$，考虑 $u(x,t)=A e^{-n^2t}\sin(nx)$。扩散系数归一化为一。

\paragraph{命题：正弦模式满足热方程}

上述模式满足 $u_t=u_{xx}$，且初值为 $u(x,0)=A\sin(nx)$。有限个此类模式相加仍满足该方程。

\begin{proof}

对时间求导得到 $-n^2A e^{-n^2t}\sin(nx)$；对空间求两次导数得到同一个表达式。整数频率保证周期边界匹配。

对有限和可以逐项求导，故线性组合仍满足方程。推广到无限级数时，逐项求导需要额外收敛条件；有限叠加的验证不自动证明任意初值问题的解存在。

\end{proof}
